\documentclass{beamer}
%[aspectratio=169]   \usepackage[czech]{babel}
\usepackage{apo-lecture-cz}
\usepackage{pdfpages}
\usepackage{pdfcomment}
\usepackage{listings}
\usepackage{array,multirow}

\subtitle{Lekce 03. Central Processing Unit (CPU)}
\author{Pavel Píša \phantom{xxxxxxx} Petr Štěpán \\ \small\texttt{pisa@fel.cvut.cz}\phantom{xxxx}\small\texttt{stepan@fel.cvut.cz}}
\begin{document}

\maketitle

\section{Procesor}

\input{common/cpu/tex/siglecycle-riscv-goal-cz.tex}

\input{common/simulator/tex/qtrvsim-introduction-cz.tex}

\input{common/cpu/tex/cpu-cycle-basic-and-asm-cz.tex}

\section{Kódování instrukcí}

\input{common/cpu/tex/register-gpr-fpr-pc-riscv-cz.tex}

\input{common/cpu/tex/instruction-encoding-mips-cz.tex}

\input{common/cpu/tex/instruction-encoding-riscv-cz.tex}

\section{Konstrukce procesoru}

\input{common/cpu/tex/cpu-cycle-pc-only-cz.tex}

\input{common/cpu/tex/cpu-building-blocks-cz.tex}

\input{common/cpu/tex/cpu-design-lw-sw-addi-alu-branch-cz.tex}

\input{common/cpu/tex/cpu-single-cycle-speed-cz.tex}

\input{common/cpu/tex/cpu-control-unit-riscv-cz.tex}


\begin{frame}
\frametitle{Bonusový bod}

Které bity instrukce rozhodují o tom, jestli se bude ukládat do paměti, nebo počítat s dvěmi registry a výsledek se uloží do třetího registru?
\begin{itemize}
\item[A] opcode
\item[B] fnct3,7
\item[C] rs1, rs2
\item[D] rd
\end{itemize}
\end{frame}

\section{Paměť úvod}

\input{common/memory/tex/memory-introduction-cz.tex}

\section{Technologie polovodičové paměťi}

\input{common/memory/tex/memory-kinds-cz.tex}

\input{common/memory/tex/memory-dimm-sdram-cz.tex}

\input{common/memory/tex/memory-gap-cz.tex}

\section{Zrychlení přístupu k datům, hierarchie, vyrovnávací paměť}

\input{common/memory/tex/memory-hierarchy-cz.tex}

\input{common/memory/tex/memory-cache-basis-cz.tex}

\input{common/simulator/tex/qtrvsim-cache-examples-cz.tex}

\input{common/memory/tex/memory-cache-parameter-impact-cz.tex}

\begin{frame}
\frametitle{Odkazy a literatura}

\begin{itemize}
\item Ulrich Drepper: \href{https://lwn.net/Articles/250967/}{What every programmer should know about memory}
\item Agner Fog: \href{https://www.agner.org/optimize/}{Software optimization resources. C++ and assembly}
\item \url{https://www.7-cpu.com/cpu/Haswell.html}
\item \url{https://www.7-cpu.com/cpu/Skylake.html}
\item WikiChips: \href{https://en.wikichip.org/wiki/amd/microarchitectures/zen}{Zen - Microarchitectures - AMD}

\end{itemize}

\end{frame}

\end{document}

